\documentclass[a4paper]{article}

\usepackage{graphicx}
\usepackage{amsmath,amssymb}
\usepackage{minted}
\usepackage{multirow}
\usepackage{float}
\usepackage{algorithm}
\usepackage{algpseudocode}
\usepackage{todonotes}
\usepackage{forest}
\usepackage{xstring}
\usepackage{xcolor}
\usepackage[most]{tcolorbox}
\usepackage{xparse}
\usepackage{natbib}

\title{Canonical range bookmark rewrite}
\date{2026-09-17}
\author{Mohammad Rahmani}

\begin{document}
\maketitle

\section{Feature}
Bookmarks are now canonical half-open character ranges rather than line
markers.  Source and Visual use the same range model, and the virtualized large
Source editor passes its local edit offsets into that same canonical tracker
instead of maintaining a second bookmark mechanism.

A selected range can be added from the Source or Visual context menu.  A
right-click inside an existing bookmarked range offers removal instead.  The
Bookmarks menu exposes the same contextual Add/Remove command.  Bookmarked text
is presented with a yellow background in Source and Visual without modifying
the LaTeX source.

Persisted bookmarks store a stable id, start and end character offsets, the
selected text snapshot, and an optional name.  Existing selection bookmarks are
migrated once; old line-only bookmarks are intentionally omitted from the new
range-only model.  Manage Bookmarks rebuilds from the live repository on every
refresh and provides navigation, single removal, and removal of all bookmarks.

\end{document}
